\section{Выбор метода решения и разработка вычислительного алгоритма}
\label{sec:algorithm}

\subsection{Обоснование выбора}

Прямой табличный симплекс-метод выбран для учебной реализации, поскольку
текущий базис, ведущий элемент и критерий остановки видны в таблице.
Графический метод нагляден для двух переменных, а перебор базисов требует
проверки всех сочетаний столбцов. Симплекс-метод исследует выбранные
базисы, но в худшем случае число шагов может быть экспоненциальным.
Методы внутренней точки подходят для крупных задач. Табличный метод
проще проверить вручную, хотя на каждом шаге обновляется вся плотная
таблица~\cite{vanderbei2020}.

\subsection{Приведение к задаче с начальным базисом}

Пусть $c$ обозначает исходный вектор цели. Внутренняя задача всегда
решается на максимум. Для максимизации берут $d=c$, для минимизации
берут $d=-c$. Если исходные строки сильно различаются по масштабу,
на обновления таблицы влияют различия в единицах измерения. Поэтому
для каждой строки вводится положительный масштаб $\lambda_i$.
Для ненулевой строки выбирается наибольшая степень двойки, которая
не превосходит максимума модулей коэффициентов и правой части.
Для полностью нулевой строки масштаб равен единице. Деление
на степень двойки сохраняет значащие биты, пока результат помещается
в нормализованный диапазон числа с плавающей точкой. Обозначим
$\widehat A_{i\cdot}=A_{i\cdot}/\lambda_i$ и
$\widehat b_i=b_i/\lambda_i$.

Для каждой строки вводится переменная запаса $s_i\geq 0$.
Если $b_i\geq 0$, базисной переменной можно взять $s_i$.
Если $b_i<0$, строка умножается на $-1$ и получает собственную
искусственную переменную $u_i\geq 0$. Система принимает вид
\begin{equation}
  \begin{cases}
    \widehat A_{i\cdot}x+s_i=\widehat b_i, & b_i\geq 0,\\
    -\widehat A_{i\cdot}x-s_i+u_i=-\widehat b_i, & b_i<0.
  \end{cases}
  \label{eq:simplex-augmented-system}
\end{equation}
В системе~\eqref{eq:simplex-augmented-system} исходный допустимый
базис расширенной задачи состоит из столбцов $s_i$ в первой группе
строк и $u_i$ во второй. При нулевых небазисных переменных значения
базисных переменных равны неотрицательным правым частям.

На первой фазе максимизируется вспомогательная цель
\begin{equation}
  w=-\sum_{i:\,b_i<0} u_i .
  \label{eq:simplex-auxiliary-objective}
\end{equation}
Поскольку все $u_i$ неотрицательны, верхняя граница
цели~\eqref{eq:simplex-auxiliary-objective} равна нулю. Исходная система
допустима тогда и только тогда, когда этот нуль достигается. Действительно,
допустимое решение исходной задачи продолжается до допустимого решения
расширенной системы с нулевыми $u_i$. Обратно, при $w=0$ все $u_i$
равны нулю, поэтому решение удовлетворяет исходным ограничениям.
При отрицательной вспомогательной цели и положительном искусственном
остатке программа возвращает статус \texttt{infeasible}.
При округлении вычисленное значение $w$ может оказаться неотрицательным
при положительном искусственном остатке. Тогда программа проверяет точку
по исходным ограничениям без пользовательского допуска и обнуляет остаток
только после успешной проверки. Если остатка нет, переход определяется
нулевыми искусственными переменными.

Нулевая искусственная переменная может остаться базисной.
Для её вывода выбирается первый ненулевой коэффициент при исходной
переменной или переменной запаса и выполняется ведущий шаг.
Нулевая правая часть при этом сохраняется, поэтому новые базисные
значения остаются допустимыми. Если все коэффициенты при исходных
переменных и переменных запаса нулевые, после исключения искусственных столбцов эта
строка представляет равенство $0=0$ и удаляется. Затем искусственные
столбцы исключаются, а строка цели $d^\mathsf{T}x$ приводится к
текущему базису. При отсутствии отрицательных правых частей
программа сразу переходит к этой второй фазе.

\subsection{Выбор ведущего элемента}

Обозначим через $y$ вектор всех переменных текущей фазы, через $B$ и
$N$ множества индексов базисных и небазисных переменных. В канонической
таблице строки ограничений и строка цели имеют вид
\begin{equation}
  y_{B_i}+\sum_{j\in N}t_{ij}y_j=\beta_i,\qquad
  z-\sum_{j\in N}r_jy_j=z_0,\qquad \beta_i\geq 0 .
  \label{eq:simplex-dictionary}
\end{equation}
Здесь $t_{ij}$ --- преобразованные коэффициенты, $\beta_i$ ---
текущие значения базисных переменных, а $r_j$ --- приведённые стоимости.
Из~\eqref{eq:simplex-dictionary} видно, что при $r_j>0$ увеличение $y_j$
может улучшить цель. При $r_j\leq 0$ для всех небазисных переменных
текущий базис оптимален. В программе строка цели хранит коэффициенты
$-r_j$, поэтому улучшающему столбцу соответствует отрицательный элемент.

При нескольких улучшающих столбцах выбирается столбец $k$ с наименьшим
индексом. Чтобы сохранить $\beta_i\geq 0$ при увеличении $y_k$,
из строк с $t_{ik}>0$ выбирается строка с минимальным отношением
\begin{equation}
  \theta = \min_{i:\,t_{ik}>0}\frac{\beta_i}{t_{ik}} .
  \label{eq:simplex-ratio}
\end{equation}
При равенстве отношений~\eqref{eq:simplex-ratio} выбирается наименьший
индекс выходящей базисной переменной. Эта пара правил известна
как правило Бланда. В точной арифметике она исключает зацикливание,
в том числе когда $\theta=0$ и шаг меняет базис без изменения
точки и значения цели~\cite{bland1977}.

Ведущая строка делится на $t_{ik}$. Из остальных строк, включая строку
цели, вычитаются её кратные, чтобы обнулить их элементы в столбце $k$.
Если улучшающий столбец есть, а строк с положительным элементом
в этом столбце нет, увеличение $y_k$ сохраняет допустимость без
верхней границы шага. Во второй фазе это признак статуса
\texttt{unbounded}. Для вспомогательной цели такой результат противоречит
её верхней границе и трактуется как численный сбой.

\subsection{Порядок вычислений и численные ограничения}

Псевдокод решения приведён ниже.
\begin{enumerate}
  \item Проверить входные данные, выбрать знак цели и построить
        таблицу~\eqref{eq:simplex-augmented-system}.
  \item Если имеются искусственные столбцы, решить вспомогательную
        задачу, проверить допустимость и вывести искусственные переменные
        из базиса.
  \item Исключить искусственные столбцы и записать исходную цель
        в канонической форме относительно найденного базиса.
  \item Выбирать ведущий столбец и строку, пока не найден оптимум,
        не доказана неограниченность или не исчерпан лимит ведущих шагов.
  \item При оптимуме извлечь исходные переменные, проверить допустимость
        и заново вычислить значение $c^\mathsf{T}x$ по исходным
        коэффициентам.
\end{enumerate}

Лимит относится к завершённым ведущим шагам обеих фаз, включая вывод
искусственных переменных из базиса. Проверка признаков остановки
выполняется до попытки следующего шага. Например, если допустимое
решение уже оптимально, нулевой лимит не мешает вернуть результат.
Пустая система ограничений обрабатывается теми же правилами.
Улучшающее направление в этом случае приводит к неограниченности.

Таблица хранится в типе \texttt{long double}. Знаки ведущих элементов
и приведённых стоимостей проверяются без абсолютного порога. Это
позволяет отличить малый ненулевой коэффициент от нуля, например
$10^{-12}$ в ограничении $10^{-12}x_1\leq 1$. При вычитании с участием
строки, уже затронутой округлением, могут обнуляться остатки
в пределах восьми машинных эпсилонов относительно элемента до
обновления. Малые разности при обновлении точных строк сохраняются.
Точные равенства отношений в используемой арифметике
обрабатываются правилом Бланда. Лимит итераций ограничивает вычисления
в случаях, когда округление меняет поведение теоретического алгоритма.

Пусть $\varepsilon$ --- допустимая относительная невязка, по умолчанию
$10^{-9}$. Для возвращаемого решения проверяется условие
\begin{equation}
  \max\{0,\,A_{i\cdot}x-b_i\}
  \leq \varepsilon\left(|b_i|+\sum_j |a_{ij}x_j|\right)
  \quad (i=1,\ldots,m).
  \label{eq:simplex-relative-residual}
\end{equation}
Проверка~\eqref{eq:simplex-relative-residual} проводится на исходных
строках, делённых на $\lambda_i$. При потере допустимости в таблице,
недопустимой итоговой невязке или непредставимом численном результате
решатель возвращает \texttt{numerical\_failure}. Численные проверки
снижают риск выдачи ошибочного решения, однако сильно
плохо обусловленные задачи могут требовать более точной арифметики.

Если $q$ --- число искусственных столбцов, таблица содержит
$(m+1)(n+m+q+1)$ элементов на первой фазе. Обновление одного ведущего
элемента требует $O((m+1)(n+m+q+1))$ операций, не считая поиска
ведущей строки и столбца, стоимость которого линейна по их длине.

Для примера из подраздела~\ref{sec:worked-example} достаточно второй фазы.
В базис последовательно входят $x_1$ и $x_2$, выходят переменные запаса
ограничений $x_1\leq 2$ и $x_1+x_2\leq 4$. За два шага получен оптимум
$(2,2)$ со значением цели $10$.
